% Generated draft: requires mathematical and source review.
\documentclass[11pt]{article}
\usepackage[margin=25mm]{geometry}
\usepackage{fontspec}
\setmainfont{Latin Modern Roman}
\usepackage{amsmath,amssymb,mathtools,unicode-math}
\setmathfont{Latin Modern Math}
\usepackage{enumitem,needspace,xurl,hyperref}
\setlength{\parindent}{0pt}
\setlength{\parskip}{5pt}
\setlength{\emergencystretch}{3em}
\newcommand{\N}{\mathbb{N}}
\newcommand{\Z}{\mathbb{Z}}
\newcommand{\Q}{\mathbb{Q}}
\newcommand{\R}{\mathbb{R}}
\newcommand{\C}{\mathbb{C}}
\newcommand{\F}{\mathbb{F}}
\newcommand{\A}{\mathbb{A}}
\newcommand{\PP}{\mathbb{P}}
\newcommand{\E}{\mathbb{E}}
\newcommand{\eps}{\varepsilon}
\newcommand{\dd}{\,\mathrm{d}}
\begin{document}
\section*{UnsolvedMath ID: 701448}
\textbf{Source problem code:} catalogue entry
\subsection{Definitions and mathematical statement}
A variety of groups \(\mathcal V\) is a class of groups defined by group identities. A group \(F\in\mathcal V\) is relatively free on a basis \(B\subseteq F\) if \(B\) generates \(F\) and every map from \(B\) to any group \(H\in\mathcal V\) extends uniquely to a homomorphism \(F\to H\). Its rank is \(|B|\); the question concerns infinite rank. A group \(G\) is torsion-free if, for every \(g\in G\) and integer \(n\geq 1\), \(g^n=1\) implies \(g=1\).

The supplied record does not define \(\aleph_1\)-homogeneity. One possible precise interpretation is the following automorphism-extension property in the first-order language of groups \(\{\cdot,{}^{-1},1\}\): whenever \(I\) is finite or countably infinite and tuples \((a_i)_{i\in I}\) and \((b_i)_{i\in I}\) in \(G\) have the same complete type over the empty set, there is \(\alpha\in\operatorname{Aut}(G)\) satisfying \(\alpha(a_i)=b_i\) for every \(i\in I\). Having the same complete type means satisfying the same parameter-free first-order formulas, each involving finitely many of the indexed variables. This interpretation is conditional: the source's intended convention requires confirmation, particularly whether homogeneity means extension to an automorphism or only extension of partial elementary maps by individual elements.

Under the displayed automorphism-extension interpretation, the question is: for every variety \(\mathcal V\) and every torsion-free relatively free group \(F\in\mathcal V\) of infinite rank, must there exist a finite or countably infinite index set \(I\) and tuples \((a_i)_{i\in I},(b_i)_{i\in I}\) such that
\[
\operatorname{tp}^{F}((a_i)_{i\in I}/\varnothing)
=
\operatorname{tp}^{F}((b_i)_{i\in I}/\varnothing),
\qquad
\nexists\alpha\in\operatorname{Aut}(F)\ \forall i\in I\;\alpha(a_i)=b_i?
\]
Thus the intended assertion is that every such group fails to be \(\aleph_1\)-homogeneous, subject to resolving the terminology above.

The supplied source reports an affirmative answer for varieties in which all free groups are residually finite, attributed to O. Belegradek, \emph{Arch. Math. Logic} \textbf{51} (2012), 781--787. Here a group \(H\) is residually finite if, for each \(h\neq 1\), there exist a finite group \(Q\) and a homomorphism \(\varphi:H\to Q\) with \(\varphi(h)\neq 1\). The source's qualification concerns all relatively free groups in the variety; the torsion-free and infinite-rank hypotheses on the groups in the question remain in force. This special-case assertion is supplied attribution, not an independently verified result here.
\subsection{Short English statement}
The question asks whether every torsion-free relatively free group of infinite rank fails to be \(\aleph_1\)-homogeneous. The source reports an affirmative special case when all free groups in the variety are residually finite. The exact meaning of \(\aleph_1\)-homogeneous is not supplied; the mathematical formulation above uses an explicitly conditional automorphism-extension interpretation.
\subsection{Sources}
\begin{itemize}
\item Local UnsolvedMath registry, supplied complete source record: permanent ID \(701448\), exact source code \texttt{catalogue entry}, title \emph{Kourovka Notebook Problem 18.19}, upstream revision \texttt{d371ee917ac9647fbb06cd6fc529deec6f3dda82}. The registry labels the problem \texttt{open}; this is a supplied status, not an independently verified literature assessment.
\item \emph{The Kourovka Notebook (Unsolved Problems in Group Theory)}, original problem 18.19, proposed by O. V. Belegradek in 2014. The supplied background cites the 18th issue (2014) and the 21st-issue edition (2026), identified as \texttt{arXiv:1401.0300v48}. The registry describes the Issues 1--21 collection as edited by E. I. Khukhro and V. D. Mazurov. These bibliographic details are supplied attribution and have not been independently checked.
\item O. Belegradek, \emph{Arch. Math. Logic} \textbf{51} (2012), 781--787, cited in the supplied statement for the residually finite special case. No article title or direct link is supplied.
\end{itemize}
\end{document}
